\begin{frame}{Index 1 of 4: principle, slide, prompt}
\scriptsize
\begin{tabular}{@{}r>{\raggedright\arraybackslash}p{5.0cm}>{\raggedright\arraybackslash}p{4.9cm}p{1.2cm}@{}}
\textbf{Slide} & \textbf{Title} & \textbf{Principle} & \textbf{Prompt}\\\midrule
2 & Today's question & EQUIVALENCE IS A UNIVERSAL CLAIM & \\
3 & A smoke alarm for a strong claim & CHECK THE WORK & \\
4 & The confident AI claim & PREDICT BEFORE RUNNING & \\
5 & Words before symbols & DEFINE THE MODEL & \\
6 & From expression to circuit & MODEL VS PROCEDURE & \\
7 & A truth table checks all assignments & EXHAUSTIVE CHECK & \\
8 & One row defeats the universal claim & COUNTEREXAMPLE & \\
9 & What counts as proof here? & SOURCES, RULES, WORK, CHECK & \\
10 & The correct simplification from the activity & COMPARE CLAIMS & \\
\end{tabular}
\note{Use this index to find the slide, principle, and prompt file.}
\end{frame}
\begin{frame}{Index 2 of 4: principle, slide, prompt}
\scriptsize
\begin{tabular}{@{}r>{\raggedright\arraybackslash}p{5.0cm}>{\raggedright\arraybackslash}p{4.9cm}p{1.2cm}@{}}
\textbf{Slide} & \textbf{Title} & \textbf{Principle} & \textbf{Prompt}\\\midrule
11 & Expression and circuit agree & MODEL VS PROCEDURE & \\
12 & Prompt 01: small steps first & PREDICT, RUN, EXPLAIN & 01a\\
13 & Prompt 02: enumerate every assignment & CHECK THE WORK & 02a\\
14 & Prompt 03a: require the exact counterexample & FAILURE MUST BE SPECIFIC & 03a\\
15 & Three ways a check can speak & EVIDENCE HAS A SHAPE & \\
16 & Satisfiable means true somewhere & WITNESS VS TAUTOLOGY & \\
17 & Tautology means true everywhere & UNIVERSAL CHECK & \\
18 & Prompt 03b: do not overclaim from a witness & NAME THE LIMIT & 03b\\
19 & Boolean checks have a boundary & MODEL LIMIT & \\
\end{tabular}
\note{Use this index to find the slide, principle, and prompt file.}
\end{frame}
\begin{frame}{Index 3 of 4: principle, slide, prompt}
\scriptsize
\begin{tabular}{@{}r>{\raggedright\arraybackslash}p{5.0cm}>{\raggedright\arraybackslash}p{4.9cm}p{1.2cm}@{}}
\textbf{Slide} & \textbf{Title} & \textbf{Principle} & \textbf{Prompt}\\\midrule
20 & From Boolean outputs to states & MODEL THE NEXT STEP & \\
21 & A complete two-state badge model & EVERY TRANSITION IS EXPLICIT & \\
22 & The supplied trace & TRACE THE EVIDENCE & \\
23 & Prompt 04: print the trace & PREDICT, RUN, EXPLAIN & 04a\\
24 & An incomplete model can look plausible & MISSING RULES ARE EVIDENCE & \\
25 & Prompt 05: name the exact gap & VALIDATE THE MODEL & 05a\\
26 & Tuesday activity: defend one claim & INDIVIDUAL CHECK & \\
27 & Peer and AI-fluency check & CRITIQUE THE EVIDENCE & \\
28 & What a state model cannot establish & MODEL LIMIT & \\
\end{tabular}
\note{Use this index to find the slide, principle, and prompt file.}
\end{frame}
\begin{frame}{Index 4 of 4: principle, slide, prompt}
\scriptsize
\begin{tabular}{@{}r>{\raggedright\arraybackslash}p{5.0cm}>{\raggedright\arraybackslash}p{4.9cm}p{1.2cm}@{}}
\textbf{Slide} & \textbf{Title} & \textbf{Principle} & \textbf{Prompt}\\\midrule
29 & Exit sentence & BOUNDED CONCLUSION & \\
\end{tabular}
\note{Use this index to find the slide, principle, and prompt file.}
\end{frame}
